\documentclass[11pt]{article}

\usepackage[margin=1in]{geometry}
\usepackage[utf8]{inputenc}
\usepackage[T1]{fontenc}
\usepackage{lmodern}
\usepackage{amsmath,amssymb,amsthm}
\usepackage{booktabs}
\usepackage{microtype}
\usepackage[numbers,sort&compress]{natbib}
\usepackage{xcolor}
\usepackage[colorlinks=true,linkcolor=blue!50!black,citecolor=blue!50!black,urlcolor=blue!50!black]{hyperref}

\newcommand{\draftnote}[1]{\textcolor{red!70!black}{[\textbf{Draft:} #1]}}

\title{The Sign or a Point's Place:\\What T11's Reading of $S_3$ Will Tell Apart\\[0.5em]
\large Research note, 26 October 2026: not for circulation}

\author{
  David Barkhuizen\\
  Independent researcher
  \and
  Claude (Opus 5.5)\\
  Anthropic
}

\date{}

\begin{document}

\maketitle

\draftnote{Written at the first author's request on 26 October 2026, while T11's batch
(\texttt{indrajala-ml}'s \texttt{docs/\allowbreak workplans/\allowbreak paper-workplan.md}, D20) was
running. \textbf{None of T11's data is used here}: its runs are neither recertified nor long-tested,
and its reading follows D20. The algebra in Section~\ref{sec:algebra} was computed on 26 October from
\texttt{window\_tasks} at \texttt{indrajala-ml} 2900990. T4's counts are its registered findings
(``Findings: T4'' and T4's second reading on T2/T5's networks, in the same workplan).}

\section{The question}

T11 trains $S_3$'s, $S_4$'s and move-to-front-4's word problems with the same cell as T5. The only
change is that each is given, as two extra outputs, a pair of the machine's approximations that
together determine its state. For $S_3$ the pair is \emph{the sign} (even or odd) and \emph{a point's
place} (where item~0 now stands). Our networks have so far never been found holding the sign as one
of their factors, although the cascade decompositions put it first~\citep{krohn-rhodes-1965-algebraic-theory-machines}.
This note says why that matters, what the literature does and does not explain, and what T4, read
on T11's certified networks, can tell apart.

\section{\texorpdfstring{$S_3$}{S3}'s approximations}\label{sec:algebra}

The machine \texttt{word-problem-s3} has 6 states, the orderings of three items. Its tokens are the
identity, the two star transpositions (positions 1 and $i$) and the 3-cycle, and its output is the
item at position~1. An \emph{approximation} is a congruence's quotient: a partition of the states
that every token respects~\citep{egri-nagy-2024-covering-lemma-transformation-semigroups}. A
\emph{cover} is a set of approximations that together determine the state.
\texttt{principal\_congruences} finds four nontrivial ones:

\begin{center}
\begin{tabular}{@{}llll@{}}
\toprule
name & blocks & quotient monoid & reading \\
\midrule
\texttt{parity} & 2 & $\mathbb{Z}_2$ (2) & the sign; $S_3 / A_3$ \\
\texttt{places \{k\}}, $k = 0, 1, 2$ & 3 & $S_3$ (6) & item $k$'s position; a point stabiliser's cosets \\
\bottomrule
\end{tabular}
\end{center}

Three facts follow, each checked by enumeration:

\begin{enumerate}
\item \textbf{Any two of the four form a cover.} Their meet is the identity, so there are six covering
  pairs: the sign with any of the three places, and any two places. D20's rule (the fewest blocks in
  total, $2 + 3 = 5$ against $3 + 3 = 6$) picks the sign and \texttt{places \{0\}}. Three covers of
  $S_3$ never use the sign.
\item \textbf{The sign is independent of everything a single step shows.} On the uniform distribution
  of states, each pair (sign, \texttt{places \{k\}}) occurs once, and so does each pair (sign,
  output). The uniform distribution is the stationary one: the tokens generate $S_3$ and include the
  identity. So neither one point's place nor one step's answer says anything about the sign. The sign
  matters to the answers only through the dynamics, as the one bit that turns one place into the
  whole state.
\item \textbf{The sign is the simplest approximation and the least informative one.} Its quotient is a
  two-state counter, the abelian top of $S_3$'s composition series. A point's place has a quotient
  monoid of the whole of $S_3$, but it holds $\log_2 3$ of the state's $\log_2 6$ bits.
\end{enumerate}

$S_4$ behaves the same way. Its sign is independent of each \texttt{places \{k\}} and of the output
(each joint cell holds 3 of the 24 states). T11's second $S_4$ output, the pairings ($S_4 / V_4 \cong
S_3$, 6 blocks), \emph{determines} $S_4$'s sign. So $S_4$'s given pair contains the sign implicitly,
and $S_3$'s names it outright.

\section{What our networks did unasked}

T4 read every certified network trained on the answers alone. On the word problems, \textbf{no
factor held \texttt{parity}}, and every congruence a factor determined short of the whole state was a
stabiliser's (\texttt{places}). This was T4b, registered and met in both readings. $S_3$'s one
certified network in the first reading (seed 2) was a cover by the places of items 1 and 2. So was
probe 4a$'''$'s seed 7. In T4's second reading, $S_4$'s one named congruence was \texttt{places \{3\}}. K8's
exploratory reading of 8 $S_3$ networks said the same: through a point's stabiliser, never the sign.
The evidence is consistent but thin: a handful of certified $S_3$ networks, because T5 certified
$S_3$ once in 20.

\section{What the literature explains, and what it doesn't}

\paragraph{Parity is hard only when it is unknown.} It is tempting to explain the missing sign by
parity's hardness for gradient descent. Shalev-Shwartz, Shamir and Shammah show that gradients
carry exponentially little information about a \emph{random} parity, whatever the model
class~\citep[Section~2]{shalev-shwartz-2017-failures-gradient-based-deep-learning}. They also say the
problem ``is not difficult if $v^*$ is known and fixed in advance'' (a full parity is learnable by an
LSTM). $S_3$'s sign is a fixed running parity, one bit flipped by two known tokens: a one-bit
recurrent state learns it. \textbf{The hardness result does not apply, so it does not explain our
networks.}

\paragraph{Transformers often take the sign first.} Li, Guo and Andreas trained transformers on
$S_3$ and $S_5$ word problems~\citep{li-2025-how-language-models-track-state}. They found two
mechanisms: an associative scan, and a ``parity-associative'' one that first computes the sign as a
shallow heuristic and then resolves the rest. Which mechanism forms depends on architecture and
training. An auxiliary parity loss induced the parity-associative mechanism consistently, and that
mechanism generalised worse to longer sequences. Our networks show the opposite preference: they
never formed the sign unasked. And T11's extra output is the analogue of Li et al.'s parity loss.

\paragraph{What remains: hypotheses, not findings.}
\begin{itemize}
\item \textbf{H1, no pressure toward it.} Fact 2 means no partial computation short of the state
  correlates with the sign. A network that holds one point's place lowers its loss just as much by
  adding a second place as by adding the sign: either completes the state. The answers alone don't
  favour the sign.
\item \textbf{H2, a factor of the same kind.} A second place is the same computation as the first,
  a point moved by permutations, on another item. Factors of one form may be easier to grow side by
  side than a parity counter beside a place tracker.
\item \textbf{H3, the architecture's bias.} The window recurrence's factors read the token through
  a window and are snapped to codes. That may favour factors whose blocks change on most tokens:
  \texttt{places} moves on most tokens, while the sign changes only on the transpositions. We
  haven't measured this.
\end{itemize}
None of these is tested. T11 bears on H1 directly: asking for the sign supplies the pressure H1
says is missing.

\section{What the reading can tell apart}

D20 already schedules T4 on every T11 network that certifies (``do its factors hold the given pair,
or others?''). For $S_3$ there are three outcomes, and each means something different for the paper:

\begin{enumerate}
\item \textbf{A factor holds \texttt{parity}} (the sign as one factor, the other a place). The given
  output decided the decomposition: K4's mechanism in its strong form. A factor never seen unasked
  forms once asked for, as in Li et al.'s parity loss. K8 then reads ``never the sign
  \emph{unasked}'', and H1 gains support.
\item \textbf{The factors are two places, and the sign is only read out.} The network certifies
  through a cover it prefers, and computes the asked sign from the whole state. T11a can still stand
  on its rates, but the paper may say ``given its approximations, it learns $S_3$'', not ``it learns
  the given approximation''. On $S_3$, K4's claim is then about rates, not mechanism. H2 or H3 gains
  support.
\item \textbf{Few certify, and the failures hold one place and no sign.} The missing bit is the
  bottleneck even when asked for. That diagnoses $S_3$'s low rate in T5 too, and points the
  post-paper work at forming a parity factor beside a place.
\end{enumerate}

Division's precedent favours the first outcome. There, every nontrivial factor T4 read held exactly
one line, and always a line trained as an output: the sum arm's networks held $A + B$, never $A - B$.
On division, the outputs decide which approximations the factors hold. Whether that carries over to a
factor the networks avoid unasked is exactly what $S_3$ tests.

\paragraph{Beside \texorpdfstring{$S_3$}{S3}.} Move-to-front-4's given factor, \texttt{order \{0, 1, 3\}}, is of the kind
its networks formed unasked (T4: \texttt{order} congruences in 5 of 8 factors, \texttt{\{0, 1, 3\}}
once). So asking for it agrees with the networks' preference rather than overriding it. $S_4$'s
pairings factor sits between the two: it contains the sign, but as part of an $S_3$ quotient, not as a
lone bit. The three machines thus range from agreeing with the preference to overriding it. A
difference in mechanism across them is more informative than one pooled rate.

\section{To register before the reading}

These expectations are proposed here and have to be written into the paper workplan before T4 runs on
T11's networks:
\begin{itemize}
\item \textbf{$S_3$:} the share of certified networks with a factor that holds \texttt{parity} (and no
  finer \texttt{places}), and the share whose factors are two \texttt{places}. Expected (from
  division's precedent): at least half hold \texttt{parity}. Refuted if none does.
\item \textbf{$S_4$:} the share of factors that hold the pairings, or an approximation that determines
  the sign. \textbf{Move-to-front-4:} the share that hold \texttt{order \{0, 1, 3\}}, against T4's
  networks trained on the answers alone.
\item Reported per machine, never pooled, and read only after every certified network has passed its
  long tests (D13).
\end{itemize}

\bibliographystyle{plainnat}
\bibliography{library,extra}

\end{document}
